\section{Introducción}

\subsection{El Problema}

Cuando arrancamos con este proyecto, la verdad es que no teníamos ni idea de qué era una arquitectura. Como nos pasa a muchos que estamos empezando, lo que hicimos fue tirarnos de cabeza a escribir código como locos, sin estructura.

El resultado fue un ``desastre organizado'' (o sea, un desastre): teníamos archivos larguísimos donde se mezclaba todo (la lógica de la app con las consultas a la base de datos), código repetido por todas partes, y cada vez que queríamos añadir algo, teníamos que rezar para que no se dañara el resto.

Esto es un problema real, no solo de novatos. Como dicen Piñero González et al.~\cite{pinero2021buenas}, si desarrollas mal, el sistema se vuelve lento y los usuarios se quejan. En nuestro caso, teníamos ese código que todos decían: ``si funciona, ¡no lo toques!'', pero sabíamos que eso no iba a durar.

Además, sin darnos cuenta, estábamos usando ``antipatrones'' (lo que NO se debe hacer):

\begin{itemize}
    \item \textbf{God Object}: Clases que hacían de todo, como el ``Dios'' del proyecto.
    \item \textbf{Spaghetti Code}: Código enredado, como un plato de espagueti.
    \item \textbf{Copy-Paste Programming}: Duplicábamos la lógica en vez de reutilizarla.
    \item \textbf{Hard Coding}: Poníamos valores fijos directamente en el código, en lugar de en un archivo de configuración.
\end{itemize}

\subsection{La Motivación}

Nuestra principal motivación fue simple: queríamos sentirnos orgullosos de lo que hacíamos. No solo que la aplicación funcionara, sino que fuera mantenible, escalable y segura. Nos dimos cuenta de que si seguíamos así, en unos meses nadie (ni nosotros mismos) entendería ese código.

Según Mercado \& Zapata~\cite{mercado2019revision}, cuando aplicas buenas prácticas, no solo mejora la calidad, sino que también el equipo trabaja mejor y se avanza más rápido. Eso era justo lo que necesitábamos. Queríamos un proyecto que realmente mostrara que podemos hacer software de calidad.

\subsection{Los Objetivos}

\begin{enumerate}
    \item Mejorar el uso de buenas prácticas en el desarrollo del proyecto
    \item Implementar patrones de diseño que hagan el código más mantenible
    \item Demostrar que un proyecto bien estructurado se desarrolla más rápido y seguro
    \item Crear un sistema escalable que pueda crecer sin problemas
    \item Documentar el proceso de transformación de código caótico a código limpio
\end{enumerate}

\subsection{Importancia del Tema}

El uso de patrones de diseño y buenas prácticas no es solo una moda o algo que te piden en la universidad. Es la diferencia entre un proyecto que dura años y uno que hay que reescribir a los 6 meses. Como menciona Martin~\cite{martin2008clean} en ``Clean Code'', el código limpio no solo es más fácil de leer, sino que también es más fácil de mantener, probar y extender.

En el contexto empresarial colombiano, Espinosa \& Eraso~\cite{espinosa2024gestion} destacan que la gestión de tecnología y la aplicación de buenas prácticas en empresas de desarrollo de software son factores críticos para la competitividad y el éxito en el mercado.